% ===========================================================================
\chapter{Running commands}\label{ch:commands}

The last tool is the most powerful: running a shell command in the
workspace. With it, the agent can do what a programmer does after every
change, run the tests, and see for itself whether its change worked. That
is what turns ``a model that suggests fixes'' into ``an agent that fixes
things'': it can check its own work, and try again when it is wrong.

It is also the most dangerous tool, since a command can do anything you
can. Chapter~\ref{ch:permission} puts it behind your permission. This
chapter is about running a command well.

\section{Running a process}

Nex's \code{Process} runs another program. \code{create
Process.command("bash", ["-c", command])} sets up bash to run the
command; nothing happens until \code{start}. Before that we say where it
runs (the workspace), and that its error output should be mixed into its
ordinary output, so the model sees both, in the order they were printed.
Once it has started, we close its input. A command that waits for the
keyboard (a program asking ``Are you sure?'') would otherwise wait for
ever; with its input closed, it reads end-of-input at once and usually
gives up.

Two things can go wrong while it runs, and both would hang the agent.

\textbf{The pipe fills up.} A process's output goes through a pipe with a
small buffer. If nobody reads it, a process that prints a lot fills the
buffer and stops, waiting for room, while we wait for it to finish: each
waits for the other, for ever. So we read the output \emph{while} the
process runs, in a separate task: \code{spawn do \ldots\ end} starts its
body running alongside the rest of the program and gives back a
\code{Task}, whose \code{await} waits for the body to finish and gives its
result.

\textbf{It never finishes.} A test stuck in a loop, a server that keeps
serving. \code{process.wait(ms)} waits at most so long, and gives back the
exit code, or \code{nil} if the process is still running. Then we stop it.

Stopping it is less simple than it sounds. Killing bash is not enough: the
command it is running (and anything that started) carries on, holding the
output pipe open, and the reading task waits on. We have to stop the whole
tree of processes. Nex doesn't do that itself, but the Java platform it
runs on does, and Nex can call Java directly. First, import the Java class
at the top of \code{tools.nex}:

\begin{addbefore}{tools.nex}{before the Tool_Outcome class}{-- What running a tool gave back: some text, and whether it failed.}
import java.lang.ProcessHandle
\end{addbefore}

Now the tool. Its timeout is set when it is made.

\begin{append}{tools.nex}
class Run_Command inherit Tool
create
  make(ws: Workspace, timeout_seconds: Integer)
  require
    positive_timeout: timeout_seconds > 0
  do
    super.make(ws)
    this.timeout_seconds := timeout_seconds
  end

feature
  timeout_seconds: Integer

  name(): String do
    result := "run_command"
  end

  description(): String do
    result := "Run a shell command (bash) in the workspace folder and get its exit code "
      + "and output (standard output and standard error together). Use it to run the "
      + "project's checks or tests, or a program. Commands that wait for keyboard input "
      + "will not get any. A command running longer than "
      + count_text(timeout_seconds, "second") + " is stopped."
  end

  parameters(): Map[String, Any] do
    result := {"command": parameter("string", "The command, e.g. `nex checks.nex`.")}
  end

  required(): Array[String] do
    result := ["command"]
  end

  changes_things(): Boolean do
    result := true
  end

  describe(input: Json_Value): String do
    result := "run_command " + string_argument(input, "command", "?")
  end
\end{append}

The description tells the model the two surprises: no keyboard, and a time
limit. A model that knows a command will be stopped after two minutes
won't start a server and wait for it.

The result starts with the exit code, on a line of its own, because that
is the first thing to know about a command: did it work? Then the output,
cut to size.

\begin{continued}{tools.nex}
  run(input: Json_Value): Tool_Outcome do
    let command: String := string_argument(input, "command", "")
    if command.trim() = "" then
      result := create Tool_Outcome.error("the command is empty")
    else
      let process: Process := create Process.command("bash", ["-c", command])
      process.set_working_directory(workspace.root)
      process.set_redirect_error_to_output(true)
      process.start()
      process.close_stdin()
      -- Read the output while the command runs: a command that writes a
      -- lot would otherwise fill the pipe and wait for ever.
      let output: Task[String] := spawn do
        result := process.read_all
      end
      let code: ?Integer := process.wait(timeout_seconds * 1000)
      if ?code as exit_code then
        result := create Tool_Outcome.ok("exit code " + exit_code.to_string + "\n"
          + shortened(output.await, output_limit()))
      else
        stop_children()
        result := create Tool_Outcome.error("the command was stopped after "
          + count_text(timeout_seconds, "second") + ". Its output so far:\n"
          + shortened(output.await, output_limit()))
      end
    end
  end
\end{continued}

A command that fails, with a non-zero exit code, is not an error of the
\emph{tool}: the tool did its job, and a failing test run is exactly what
the model needs to see. Only a command that had to be stopped is reported
as an error.

\code{stop_children} is where Nex calls Java. Inside \code{with "java" do}
\ldots\ \code{end}, method calls go to Java objects. \code{ProcessHandle}
can list a process's descendants: everything it started, and everything
those started. We would like to ask for the descendants of the bash
process, but that needs \code{ProcessHandle.of}, and \code{of} is a word
Nex keeps for itself (in \code{case} \ldots\ \code{of}). There is a simple
way round, which also happens to be the right rule: the agent runs one
command at a time, so \emph{every} process below the agent belongs to the
command that ran too long. We stop them all.

\begin{continued}{tools.nex}
  private feature
    -- Stop every process this program has started. The agent runs one
    -- command at a time, so they all belong to the command that ran too
    -- long: bash, and whatever bash started.
    stop_children() do
      with "java" do
        let children := ProcessHandle.current().descendants().toList()
        let i: Integer := 0
        from until i = children.size() do
          children.get(i).destroyForcibly()
          i := i + 1
        end
      end
    end
end
\end{continued}

With all six tools written, one function gives the usual set, all working
in one workspace, with a two-minute limit on commands:

\begin{append}{tools.nex}
-- The usual set of tools, all working in `ws`.
function standard_tools(ws: Workspace): Array[Tool] do
  result := [create List_Files.make(ws), create Read_File.make(ws),
    create Search.make(ws), create Edit_File.make(ws),
    create Write_File.make(ws), create Run_Command.make(ws, 120)]
end
\end{append}

Let's run the sample project's checks, a command that writes to both
outputs and fails, and one that runs too long (with a one-second limit, so
we don't have to wait):

\begin{tryit}{try_commands.nex}
intern Tools

let ws: Workspace := fresh_workspace(user_path("sample_project"), user_path("scratch"))

function show(outcome: Tool_Outcome) do
  console.print_line((when outcome.is_error then "error: " else "" end) + outcome.text)
end

let run: Run_Command := create Run_Command.make(ws, 60)
show(run.run(json_value('{"command": "nex checks.nex"}')))
show(run.run(json_value('{"command": "echo to output; echo to errors >&2; exit 3"}')))
let impatient: Run_Command := create Run_Command.make(ws, 1)
show(impatient.run(json_value('{"command": "echo starting; sleep 60; echo never"}')))
\end{tryit}
\tryoutput{try_commands}

The checks ran in the workspace and failed, as they should: this copy has
the bug. The second command's two outputs arrive together, and the third
was stopped after one second, with what it printed before. Our agent now
has everything it needs to fix the bug. Before we let it, we should decide
who says it may.
